\section{B3D: generación 3D mediante flujo de datos}

\subsection{Motivación: eliminar el cuello de botella en la optimización}

Aunque los pasos de optimización del NeRF en CAT3D son efectivos, aún requieren decenas de segundos de iteración optimizada. Una pregunta natural es: ¿\textbf{es posible reemplazar el paso de optimización con un modelo de flujo de datos}?

B3D (Both3D) responde afirmativamente; puede generar una escena 3D completa en \textbf{menos de 7 segundos} usando una sola GPU.

\subsection{Idea central: Splatter Image}

La idea clave de B3D se basa en la Imagen Splatter: predecir los parámetros de un splat gaussiano 3D para cada píxel. La Imagen Splatter es en sí misma un método de reconstrucción pura y no puede manejar regiones no observadas. B3D introduce sobre esta base un \textbf{modelo generativo}, permitiendo que «infiera» (hallucine) el contenido de las regiones no vistas.

\begin{importantbox}{Flujo de generación en dos etapas de B3D}
\textbf{Primera etapa: generación mediante difusión}. Un modelo de difusión latente multi-vista genera simultáneamente imágenes RGB y mapas XYZ (de puntos) desde múltiples perspectivas.\\
\textbf{Segunda etapa: regresión de Gaussianos}. Una red Gaussian Head recibe las imágenes RGB y los mapas de puntos generados, regresando las propiedades restantes de cada splat gaussiano 3D correspondiente a cada píxel (opacidad y escala), produciendo una Imagen Splatter multi-vista.
\end{importantbox}

\subsection{¿Por qué se adopta un diseño en dos etapas?}

Esta elección de diseño tiene razones técnicas profundas:

\begin{itemize}
    \item Los colores RGB y las posiciones XYZ pueden obtenerse de manera confiable mediante señales de supervisión a partir de conjuntos de datos multi-vista —la información de posición se obtiene ejecutando una cadena de procesamiento (pipeline) de Estructura desde Movimiento (SfM) sobre conjuntos de datos multi-vista existentes.
    \item Pero la opacidad y la escala de los Gaussianos a nivel de píxel son difíciles de obtener con valores verdaderos confiables; por lo tanto, \textbf{no son adecuados para generarlos directamente mediante un modelo de difusión}.
    \item Investigaciones previas demuestran que, dada la información de color y posición, \textbf{regresar} opacidad y escala es mucho más fácil que \textbf{generarlos}.
\end{itemize}

\subsection{Gaussian Head multi-vista}

Los elementos clave del diseño de la red Gaussian Head:

\begin{itemize}
    \item Recibe los mapas de puntos y las imágenes RGB predichos para todas las vistas.
    \item Convierte cada vista en su correspondiente Imagen Splatter (es decir, un splat gaussiano 3D por píxel).
    \item Utiliza \textbf{atención entre vistas} (cross-view attention) para permitir que la red razona conjuntamente todas las Imágenes Splatter, asegurando consistencia 3D.
    \item Las múltiples Imágenes Splatter multi-vista combinadas representan la escena 3D completa.
\end{itemize}

\subsection{Datos de entrenamiento}

El entrenamiento de B3D requiere conjuntos de datos multi-vista a gran escala; el equipo utilizó MASt3R (una cadena de procesamiento de Estructura desde Movimiento 3D) para procesar todas las escenas de los siguientes conjuntos de datos:

\begin{itemize}
    \item CO3D
    \item RealEstate10K (R10K)
    \item MVImgNet
    \item DL3DV
\end{itemize}

Esto garantiza que todos los datos de entrenamiento tengan anotaciones XYZ densas. Combinado con conjuntos de datos de objetos sintéticos, se construyó un conjunto de entrenamiento a gran escala.

\begin{knowledgebox}{Estrategia de construcción de datos de entrenamiento}
Aunque los datos 3D en internet son escasos, los datos de imágenes multi-vista son relativamente abundantes. Ejecutando una cadena de procesamiento (pipeline) SfM sobre conjuntos de datos multi-vista, se pueden obtener automáticamente anotaciones de profundidad densa/mapas de puntos para cada imagen, proporcionando así suficientes señales de supervisión para el entrenamiento del modelo de difusión. Esta estrategia de «obtener supervisión 3D indirectamente mediante datos 2D» es un paradigma importante en la investigación actual de generación 3D.
\end{knowledgebox}

\subsection{Resumen de la sección}

B3D mejora la velocidad de generación 3D desde aproximadamente 1 minuto hasta menos de 7 segundos al introducir un Gaussian Head de flujo de datos que reemplaza el paso de optimización del NeRF. Su diseño en dos etapas (generación mediante difusión de RGB+XYZ, regresión de opacidad+escala) aprovecha con astucia las diferencias en la supervisabilidad entre diferentes propiedades.
